\section{What the Shared Region Carries}\label{sec:carries}

\subsection{Removing}\label{sec:remove}

We first remove the shared region directly to test whether it is needed for thinking. Removing it from the $\lambda=0$ adapter drops reflective markers to 0.2--10\% of their original level (14--44 per 1,000 words), whereas removing the same energy from other directions outside the shared region leaves 8--24, and scaling down the whole update by the same amount leaves 5--28 (Table~\ref{tab:6}). Removing the same energy from the top singular components of the update \citep{lee2025goldswitch} leaves 22--35\% of the original markers in pair A and 1--17\% in pairs B and C, and shortens responses by only 1.8--4.0$\times$, against 3.7--5.4$\times$ when the shared region is removed (Table~\ref{tab:6}). A length-matched cut of the original adapter keeps 11--22 markers per 1,000 words, while the adapter with the shared region removed keeps 0.07--1.5 (Table~\ref{tab:6}).

Removing the shared region raises accuracy from 10--18\% to 20--36\%, a mean of 24.7\% over the nine cells, higher than the other three edits (16.8--21.0\%). The length-matched cut averages 16.0\% and is below removal in all nine cells, although the gap is only just over 1 pp in A/Qwen and A/Llama (Table~\ref{tab:7}). Applying the same edits to the $\lambda=1$ adapter changes little, since it has almost no reflective markers to begin with, and response length changes by only 0.6--1.2$\times$, far less than at $\lambda=0$. Only in pair B, whose short answers still use a few reflective markers, do Qwen and OLMo drop to 27--40\% of their original level (Table~\ref{tab:6}).

Because removing the same energy has different effects depending on where it is removed, the region is needed for thinking. Removing the top components comes close to the shared-region effect in pairs B and C but not in pair A, so across all three pairs it is the shared region, not the large components of the update, that turns thinking off. Whether the drop in reflective markers when other directions are removed in pair A reflects less thinking is revisited with thinking features in \S\ref{sec:sae}. Since the same edits have small effects on the $\lambda=1$ adapter, the change at $\lambda=0$ is not a side effect of the edit itself. Responses also become 4--5$\times$ shorter, but reflective markers and accuracy still diverge at matched length, so the effect of removing the region is not explained by shorter responses alone; the accuracy gain likewise comes not only from shortening but from the student writing a short solution from start to finish. In short, removing the shared region, unlike a length-matched cut, turns thinking off and lets the student finish a short solution, but removal alone cannot tell whether the region is sufficient to turn thinking on.

\begin{table*}[t]
\centering
\small
% Table 6 (wide: use table*), generated from overleaf/05_what_the_shared_region_carries.md
\resizebox{\linewidth}{!}{%
\begin{tabular}{llrrrrrrrrr}
\toprule
Pair & Student & $\lambda$ = 0 & Shared removed & Length-matched cut & Other removed & Top removed (matched) & Scaled down & Tokens ($\lambda$ = 0 $\rightarrow$ shared removed) & $\lambda$ = 1 markers ($\lambda$ = 1 $\rightarrow$ shared removed) & $\lambda$ = 1 tokens ($\lambda$ = 1 $\rightarrow$ shared removed) \\
\midrule
A & Qwen2.5-0.5B & 21.8 $\pm$ 2.5 & 0.37 $\pm$ 0.13 & 13.8 $\pm$ 0.6 & 19.1 $\pm$ 2.1 & 6.0 $\pm$ 1.2 & 13.0 $\pm$ 1.1 & 5,916 $\pm$ 214 $\rightarrow$ 1,439 $\pm$ 80 & 0.48 $\pm$ 0.61 $\rightarrow$ 0.63 $\pm$ 0.81 & 811 $\pm$ 64 $\rightarrow$ 860 $\pm$ 52 \\
 & Llama-3.2-1B & 30.7 $\pm$ 1.2 & 0.18 $\pm$ 0.05 & 19.4 $\pm$ 1.2 & 17.9 $\pm$ 0.5 & 10.6 $\pm$ 2.4 & 27.6 $\pm$ 5.3 & 6,969 $\pm$ 77 $\rightarrow$ 1,659 $\pm$ 211 & 0.12 $\pm$ 0.03 $\rightarrow$ 0.49 $\pm$ 0.42 & 944 $\pm$ 72 $\rightarrow$ 856 $\pm$ 32 \\
 & OLMo-2-1B & 15.0 $\pm$ 1.5 & 1.2 $\pm$ 0.1 & 11.2 $\pm$ 0.6 & 8.8 $\pm$ 1.6 & 3.3 $\pm$ 1.0 & 10.3 $\pm$ 0.3 & 3,015 $\pm$ 86 $\rightarrow$ 732 $\pm$ 16 & 0.41 $\pm$ 0.13 $\rightarrow$ 0.37 $\pm$ 0.16 & 683 $\pm$ 23 $\rightarrow$ 558 $\pm$ 15 \\
\midrule
B & Qwen2.5-0.5B & 19.2 $\pm$ 2.2 & 0.15 $\pm$ 0.03 & 14.0 $\pm$ 0.6 & 15.2 $\pm$ 3.2 & 1.2 $\pm$ 0.8 & 4.6 $\pm$ 0.4 & 5,226 $\pm$ 82 $\rightarrow$ 966 $\pm$ 25 & 4.7 $\pm$ 2.1 $\rightarrow$ 1.8 $\pm$ 0.46 & 2,302 $\pm$ 94 $\rightarrow$ 1,839 $\pm$ 158 \\
 & Llama-3.2-1B & 44.4 $\pm$ 26.0 & 0.07 $\pm$ 0.03 & 22.1 $\pm$ 7.0 & 12.8 $\pm$ 7.1 & 0.58 $\pm$ 0.47 & 6.8 $\pm$ 1.5 & 5,095 $\pm$ 164 $\rightarrow$ 1,141 $\pm$ 80 & 2.8 $\pm$ 0.72 $\rightarrow$ 4.0 $\pm$ 2.0 & 1,765 $\pm$ 78 $\rightarrow$ 2,065 $\pm$ 3 \\
 & OLMo-2-1B & 16.8 $\pm$ 2.4 & 1.4 $\pm$ 0.2 & 14.3 $\pm$ 1.5 & 7.7 $\pm$ 1.2 & 1.7 $\pm$ 0.3 & 5.4 $\pm$ 2.6 & 2,827 $\pm$ 25 $\rightarrow$ 760 $\pm$ 24 & 5.1 $\pm$ 2.0 $\rightarrow$ 1.4 $\pm$ 0.02 & 1,207 $\pm$ 40 $\rightarrow$ 825 $\pm$ 35 \\
\midrule
C & Qwen2.5-0.5B & 16.6 $\pm$ 1.7 & 0.38 $\pm$ 0.24 & 11.7 $\pm$ 0.2 & 24.3 $\pm$ 4.0 & 2.9 $\pm$ 0.3 & 8.4 $\pm$ 1.4 & 5,231 $\pm$ 141 $\rightarrow$ 1,025 $\pm$ 67 & 0.07 $\pm$ 0.06 $\rightarrow$ 0.03 $\pm$ 0.01 & 1,399 $\pm$ 156 $\rightarrow$ 1,313 $\pm$ 91 \\
 & Llama-3.2-1B & 19.1 $\pm$ 1.1 & 0.10 $\pm$ 0.12 & 16.3 $\pm$ 0.4 & 22.2 $\pm$ 2.7 & 0.81 $\pm$ 0.29 & 13.8 $\pm$ 0.8 & 6,149 $\pm$ 103 $\rightarrow$ 1,214 $\pm$ 118 & 0.14 $\pm$ 0.14 $\rightarrow$ 0.05 $\pm$ 0.01 & 939 $\pm$ 96 $\rightarrow$ 912 $\pm$ 45 \\
 & OLMo-2-1B & 14.4 $\pm$ 1.5 & 1.4 $\pm$ 0.03 & 11.9 $\pm$ 1.1 & 12.9 $\pm$ 1.5 & 1.5 $\pm$ 0.3 & 8.8 $\pm$ 1.6 & 2,949 $\pm$ 50 $\rightarrow$ 744 $\pm$ 44 & 0.33 $\pm$ 0.22 $\rightarrow$ 0.37 $\pm$ 0.31 & 910 $\pm$ 37 $\rightarrow$ 567 $\pm$ 57 \\
\bottomrule
\end{tabular}}

\caption{Reflective-marker frequency (per 1,000 words) and response length when the shared region is removed from the $\lambda=0$ adapter. Other removed removes the same energy from other directions outside the shared region, Top removed (matched) removes the same energy from the top singular components of the update, Scaled down scales the whole update down by the same amount, and Length-matched cut cuts the original adapter at the same mean length as the adapter with the shared region removed. The last two columns are a control that applies the same edit to the $\lambda=1$ adapter. Reflective markers disappear most clearly when the shared region is removed. Values after $\pm$ are standard deviations across three seeds.}
\label{tab:6}
\end{table*}

\begin{table*}[t]
\centering
\small
% Table 7 (wide: use table*), generated from overleaf/05_what_the_shared_region_carries.md
\resizebox{\linewidth}{!}{%
\begin{tabular}{llrrrrrr}
\toprule
Pair & Student & $\lambda$ = 0 & Shared removed & Length-matched cut & Other removed & Top removed (matched) & Scaled down \\
\midrule
A & Qwen2.5-0.5B & 15.2 $\pm$ 1.4 & 20.2 $\pm$ 2.3 & 18.8 $\pm$ 2.0 & 16.9 $\pm$ 2.8 & 18.8 $\pm$ 3.2 & 18.3 $\pm$ 1.7 \\
 & Llama-3.2-1B & 11.8 $\pm$ 1.3 & 22.1 $\pm$ 1.5 & 21.0 $\pm$ 0.9 & 19.4 $\pm$ 2.2 & 17.2 $\pm$ 1.5 & 16.8 $\pm$ 0.9 \\
 & OLMo-2-1B & 10.2 $\pm$ 2.3 & 19.8 $\pm$ 0.8 & 11.2 $\pm$ 1.4 & 14.5 $\pm$ 2.6 & 21.8 $\pm$ 1.4 & 17.3 $\pm$ 0.8 \\
\midrule
B & Qwen2.5-0.5B & 17.5 $\pm$ 0.7 & 35.8 $\pm$ 0.4 & 14.7 $\pm$ 1.2 & 19.2 $\pm$ 1.9 & 25.3 $\pm$ 3.6 & 22.2 $\pm$ 1.7 \\
 & Llama-3.2-1B & 17.1 $\pm$ 0.9 & 23.7 $\pm$ 2.8 & 20.1 $\pm$ 0.8 & 21.4 $\pm$ 4.5 & 21.9 $\pm$ 0.5 & 21.9 $\pm$ 2.4 \\
 & OLMo-2-1B & 10.5 $\pm$ 2.0 & 21.4 $\pm$ 1.3 & 11.2 $\pm$ 0.1 & 14.5 $\pm$ 0.4 & 20.4 $\pm$ 1.0 & 19.2 $\pm$ 1.4 \\
\midrule
C & Qwen2.5-0.5B & 16.2 $\pm$ 1.1 & 34.5 $\pm$ 4.1 & 17.4 $\pm$ 0.7 & 16.8 $\pm$ 1.8 & 20.9 $\pm$ 1.3 & 20.7 $\pm$ 1.4 \\
 & Llama-3.2-1B & 15.0 $\pm$ 2.0 & 23.6 $\pm$ 1.5 & 19.5 $\pm$ 2.1 & 14.5 $\pm$ 1.1 & 21.2 $\pm$ 1.1 & 21.8 $\pm$ 1.2 \\
 & OLMo-2-1B & 11.7 $\pm$ 1.6 & 20.9 $\pm$ 0.8 & 10.5 $\pm$ 0.3 & 13.7 $\pm$ 0.1 & 21.0 $\pm$ 0.5 & 20.2 $\pm$ 0.6 \\
\midrule
Mean &  & 13.9 & 24.7 & 16.0 & 16.8 & 21.0 & 19.8 \\
\bottomrule
\end{tabular}}

\caption{MATH-500 accuracy (\%) of the edited $\lambda=0$ adapter. Length-matched cut cuts the original adapter at the same mean length as the adapter with the shared region removed and then has it write the answer; the other columns are as in Table~\ref{tab:6}. The adapter with the shared region removed is the most accurate and is above the length-matched cut in all nine cells. Values after $\pm$ are standard deviations across three seeds.}
\label{tab:7}
\end{table*}

\subsection{Grafting}\label{sec:graft}

We therefore graft the shared region into an adapter that only gave short answers, to see whether the region alone makes the student think. Adding only the shared-region component of the $\lambda=0$ adapter to the $\lambda=1$ adapter ($\beta=1$) gives a recovery rate, measured above the $\lambda=1$ level, of 134\% in B/OLMo, beyond the $\lambda=0$ level, but only 1--16\% in the other eight cells, while response length alone grows by 1.3--3.3$\times$ (Table~\ref{tab:8}). In B/OLMo, 87\% of responses hit the generation limit and accuracy falls to 35\% of the $\lambda=1$ adapter's, so the student may have repeated the same phrases. Grafting the same energy from other directions outside the shared region mostly changes little, but in B/Llama it raises markers more than the shared region does; grafting the whole $\lambda=0$ update at the same energy is weaker than grafting the shared region alone in every cell except A/Llama (Table~\ref{tab:8}).

Raising the grafted amount to $\beta=2$, the shared region turns on reflective markers in the early part of responses in all nine cells, at 5.7--28 per 1,000 words. This exceeds the $\lambda=0$ adapter (10--16) in seven cells and reaches 56--68\% of it in A/Llama and A/OLMo, whereas grafting from other directions yields only 0.08--3.4 (Table~\ref{tab:9}). However, adapters with the shared region grafted finish at most 1\% of responses, except A/Qwen and A/Llama (17--23\%), and their accuracy falls to 3--9\%, while adapters with other directions grafted finish 55--91\% of responses at 10--29\% accuracy (Table~\ref{tab:9}).

Cells that did not recover at $\beta=1$ turn on reflective markers at $\beta=2$, and grafting the same energy along other directions does not, so it is this region that turns thinking on. B/Llama, however, does not turn on thinking features (\S\ref{sec:sae}), so in pair B the student may only have repeated the phrases. The thinking that is turned on does not lead to finished solutions or correct answers. Adapters grafted with other directions at the same energy do finish, so the failure to finish is due to this region rather than to the grafted amount, and what is needed to finish a solution is not transferred by this region alone. A refusal direction induces refusal when merely added \citep{arditi2024refusal}, but this region, once grafted, only starts thinking without finishing the solution, so removal and grafting are not symmetric. As with the features reported to be necessary and sufficient for hesitation markers \citep{hu2026transcoder}, this region is sufficient to turn thinking on, but not to finish a solution. In short, grafting the region starts thinking but does not finish the solution: the region turns thinking on rather than holding the whole solution.

\begin{table*}[t]
\centering
\small
% Table 8 (wide: use table*), generated from overleaf/05_what_the_shared_region_carries.md
\resizebox{\linewidth}{!}{%
\begin{tabular}{llrrrrrr}
\toprule
Pair & Student & Markers: $\lambda$ = 1 & Markers: graft shared & Markers: graft other & Markers: graft whole & Tokens ($\lambda$ = 1 $\rightarrow$ graft shared) & Recovered (\%) \\
\midrule
A & Qwen2.5-0.5B & 0.48 $\pm$ 0.61 & 1.7 $\pm$ 1.4 & 0.24 $\pm$ 0.07 & 0.26 $\pm$ 0.12 & 811 $\pm$ 64 $\rightarrow$ 1,088 $\pm$ 103 & 6 \\
 & Llama-3.2-1B & 0.12 $\pm$ 0.03 & 1.1 $\pm$ 0.2 & 0.51 $\pm$ 0.28 & 1.4 $\pm$ 0.7 & 944 $\pm$ 72 $\rightarrow$ 1,969 $\pm$ 138 & 3 \\
 & OLMo-2-1B & 0.41 $\pm$ 0.13 & 0.82 $\pm$ 0.56 & 0.32 $\pm$ 0.23 & 0.45 $\pm$ 0.14 & 683 $\pm$ 23 $\rightarrow$ 1,699 $\pm$ 108 & 3 \\
\midrule
B & Qwen2.5-0.5B & 4.7 $\pm$ 2.1 & 6.9 $\pm$ 1.3 & 2.9 $\pm$ 0.6 & 4.0 $\pm$ 0.1 & 2,302 $\pm$ 94 $\rightarrow$ 5,795 $\pm$ 147 & 16 \\
 & Llama-3.2-1B & 2.8 $\pm$ 0.7 & 5.9 $\pm$ 1.9 & 7.2 $\pm$ 6.2 & 4.4 $\pm$ 1.7 & 1,765 $\pm$ 78 $\rightarrow$ 5,790 $\pm$ 470 & 7 \\
 & OLMo-2-1B & 5.1 $\pm$ 2.0 & 20.8 $\pm$ 8.6 & 4.6 $\pm$ 3.4 & 18.8 $\pm$ 3.9 & 1,207 $\pm$ 40 $\rightarrow$ 3,467 $\pm$ 30 & 134 \\
\midrule
C & Qwen2.5-0.5B & 0.07 $\pm$ 0.06 & 1.7 $\pm$ 0.5 & 0.14 $\pm$ 0.11 & 1.1 $\pm$ 0.8 & 1,399 $\pm$ 156 $\rightarrow$ 4,298 $\pm$ 143 & 10 \\
 & Llama-3.2-1B & 0.14 $\pm$ 0.14 & 0.81 $\pm$ 0.41 & 0.35 $\pm$ 0.46 & 0.48 $\pm$ 0.13 & 939 $\pm$ 96 $\rightarrow$ 1,470 $\pm$ 25 & 3 \\
 & OLMo-2-1B & 0.33 $\pm$ 0.22 & 0.47 $\pm$ 0.61 & 0.29 $\pm$ 0.36 & 0.20 $\pm$ 0.21 & 910 $\pm$ 37 $\rightarrow$ 1,900 $\pm$ 173 & 1 \\
\bottomrule
\end{tabular}}

\caption{Reflective-marker frequency and response length when components of the $\lambda=0$ adapter are added to the $\lambda=1$ adapter. graft shared adds the shared-region component, graft other adds the same energy taken from outside the shared region, and graft whole adds the whole $\lambda=0$ update scaled to the same energy. Recovered is the recovery rate measured above the $\lambda=1$ level, (graft shared $-$ $\lambda=1$)/($\lambda=0$ $-$ $\lambda=1$). The table separates whether the shared region alone brings back thinking or only length. Values after $\pm$ are standard deviations across three seeds.}
\label{tab:8}
\end{table*}

\begin{table*}[t]
\centering
\small
% Table 9 (wide: use table*), generated from overleaf/05_what_the_shared_region_carries.md
\resizebox{\linewidth}{!}{%
\begin{tabular}{llrrrrrrrr}
\toprule
Pair & Student & Markers: $\lambda$ = 0 & Markers: graft shared & Markers: graft other & Markers: graft whole & Completion rate: graft shared & Completion rate: graft other & Acc. (\%): graft shared & Acc. (\%): graft other \\
\midrule
A & Qwen2.5-0.5B & 10.2 $\pm$ 0.8 & 19.6 $\pm$ 1.5 & 0.41 $\pm$ 0.11 & 12.1 $\pm$ 1.6 & 0.23 $\pm$ 0.11 & 0.91 $\pm$ 0.02 & 7.7 $\pm$ 1.3 & 28.9 $\pm$ 2.1 \\
 & Llama-3.2-1B & 12.4 $\pm$ 0.04 & 8.4 $\pm$ 0.7 & 1.0 $\pm$ 0.3 & 8.3 $\pm$ 0.1 & 0.17 $\pm$ 0.07 & 0.63 $\pm$ 0.07 & 8.5 $\pm$ 1.1 & 18.8 $\pm$ 1.3 \\
 & OLMo-2-1B & 10.1 $\pm$ 0.5 & 5.7 $\pm$ 0.6 & 0.29 $\pm$ 0.12 & 10.7 $\pm$ 0.3 & 0.00 $\pm$ 0.003 & 0.75 $\pm$ 0.03 & 5.0 $\pm$ 0.4 & 10.2 $\pm$ 0.8 \\
\midrule
B & Qwen2.5-0.5B & 15.7 $\pm$ 1.4 & 20.4 $\pm$ 4.2 & 1.1 $\pm$ 0.3 & 11.9 $\pm$ 0.6 & 0.00 $\pm$ 0.00 & 0.59 $\pm$ 0.08 & 3.5 $\pm$ 0.4 & 19.2 $\pm$ 1.8 \\
 & Llama-3.2-1B & 16.2 $\pm$ 0.2 & 18.7 $\pm$ 7.2 & 3.4 $\pm$ 0.7 & 70.4 $\pm$ 50.3 & 0.00 $\pm$ 0.001 & 0.55 $\pm$ 0.16 & 3.3 $\pm$ 1.1 & 13.2 $\pm$ 1.5 \\
 & OLMo-2-1B & 15.1 $\pm$ 0.6 & 28.0 $\pm$ 11.2 & 0.80 $\pm$ 0.19 & 59.0 $\pm$ 40.2 & 0.00 $\pm$ 0.003 & 0.74 $\pm$ 0.04 & 3.1 $\pm$ 0.4 & 9.5 $\pm$ 1.5 \\
\midrule
C & Qwen2.5-0.5B & 10.2 $\pm$ 0.3 & 17.0 $\pm$ 6.8 & 0.16 $\pm$ 0.04 & 1.7 $\pm$ 0.6 & 0.00 $\pm$ 0.00 & 0.68 $\pm$ 0.01 & 4.7 $\pm$ 0.4 & 17.8 $\pm$ 2.0 \\
 & Llama-3.2-1B & 12.5 $\pm$ 0.2 & 27.1 $\pm$ 2.8 & 0.23 $\pm$ 0.05 & 16.9 $\pm$ 0.6 & 0.00 $\pm$ 0.003 & 0.78 $\pm$ 0.02 & 4.0 $\pm$ 1.1 & 24.8 $\pm$ 1.3 \\
 & OLMo-2-1B & 10.1 $\pm$ 0.3 & 10.2 $\pm$ 1.1 & 0.08 $\pm$ 0.06 & 10.2 $\pm$ 0.6 & 0.00 $\pm$ 0.001 & 0.62 $\pm$ 0.06 & 4.1 $\pm$ 0.9 & 9.7 $\pm$ 2.0 \\
\bottomrule
\end{tabular}}

\caption{Reflective-marker frequency, share of completed responses, and accuracy when the grafted amount is raised to $\beta=2$. graft shared, graft other and graft whole are as in Table~\ref{tab:8}; frequencies are measured on the early part of each response, before repetition builds up. Thinking turns on only when the shared region is grafted, but those responses are not finished. Values after $\pm$ are standard deviations across three seeds.}
\label{tab:9}
\end{table*}

\subsection{Amplifying}\label{sec:amplify}

We therefore increase the weight of the shared region to see whether it also controls the amount of thinking. When the norm of the $\lambda=0$ adapter is kept fixed and only the weight on the shared region is raised ($\gamma = 1, 4$), reflective markers do not increase in pairs A and B or in C/Llama and C/OLMo, but fall to or stay at 26--100\% of the original; only in C/Qwen do they rise, by 1.3--1.6$\times$, and there response length and truncation rise with them (Table~\ref{tab:10}). Even moving the entire update into the shared region leaves reflective markers at 22--103\% of the original, except in C/Qwen (1.7$\times$) (Table~\ref{tab:10}).

Since markers do not increase in most cells when the region is scaled up under a fixed norm, the region's weight does not set the amount of thinking. This differs from reports that strengthening activation directions or features lengthens reasoning \citep{sheng2025planning,jayabahu2026halt,galichin2026covered}. Those works scale the activations of a model that already thinks, without a norm constraint, whereas here we only shift weight within the weight update while holding its norm fixed. Because 51--87\% of the $\lambda=0$ adapter's responses already hit the generation limit, leaving little room to grow longer, we cannot fully rule out a ceiling effect. The rise in C/Qwen comes together with longer responses and more truncation, so it is hard to read as thinking alone being amplified. In short, removing the shared region turns thinking off but amplifying it does not increase thinking, so the region is closer to an on/off switch than to a dial that sets the amount of thinking.

\begin{table*}[t]
\centering
\small
% Table 10 (fits one column or use table*), generated from overleaf/05_what_the_shared_region_carries.md
\begin{tabular}{llrrrrr}
\toprule
Pair & Student & $\lambda$ = 0 & $\gamma$ = 1 & $\gamma$ = 4 & All in region & Tokens ($\lambda$ = 0 $\rightarrow$ all in region) \\
\midrule
A & Qwen2.5-0.5B & 21.8 $\pm$ 2.5 & 21.2 $\pm$ 2.9 & 19.4 $\pm$ 2.9 & 19.1 $\pm$ 3.8 & 5,916 $\pm$ 214 $\rightarrow$ 6,089 $\pm$ 663 \\
 & Llama-3.2-1B & 30.7 $\pm$ 1.2 & 20.9 $\pm$ 0.7 & 16.2 $\pm$ 2.3 & 13.7 $\pm$ 2.0 & 6,969 $\pm$ 77 $\rightarrow$ 6,173 $\pm$ 335 \\
 & OLMo-2-1B & 15.0 $\pm$ 1.5 & 10.1 $\pm$ 0.9 & 8.7 $\pm$ 1.1 & 7.4 $\pm$ 1.2 & 3,015 $\pm$ 86 $\rightarrow$ 2,993 $\pm$ 348 \\
\midrule
B & Qwen2.5-0.5B & 19.2 $\pm$ 2.2 & 10.8 $\pm$ 1.0 & 11.3 $\pm$ 2.7 & 11.1 $\pm$ 0.8 & 5,226 $\pm$ 82 $\rightarrow$ 5,386 $\pm$ 220 \\
 & Llama-3.2-1B & 44.4 $\pm$ 26.0 & 18.8 $\pm$ 11.7 & 11.6 $\pm$ 1.5 & 9.9 $\pm$ 2.9 & 5,095 $\pm$ 164 $\rightarrow$ 5,808 $\pm$ 751 \\
 & OLMo-2-1B & 16.8 $\pm$ 2.4 & 10.8 $\pm$ 1.5 & 9.9 $\pm$ 1.0 & 9.4 $\pm$ 1.4 & 2,827 $\pm$ 25 $\rightarrow$ 2,990 $\pm$ 133 \\
\midrule
C & Qwen2.5-0.5B & 16.6 $\pm$ 1.7 & 21.7 $\pm$ 1.1 & 25.6 $\pm$ 4.2 & 27.3 $\pm$ 1.7 & 5,231 $\pm$ 141 $\rightarrow$ 7,003 $\pm$ 266 \\
 & Llama-3.2-1B & 19.1 $\pm$ 1.1 & 19.1 $\pm$ 0.7 & 18.3 $\pm$ 2.2 & 19.6 $\pm$ 4.1 & 6,149 $\pm$ 103 $\rightarrow$ 7,428 $\pm$ 99 \\
 & OLMo-2-1B & 14.4 $\pm$ 1.5 & 11.4 $\pm$ 1.6 & 12.5 $\pm$ 1.1 & 12.1 $\pm$ 1.1 & 2,949 $\pm$ 50 $\rightarrow$ 3,132 $\pm$ 217 \\
\bottomrule
\end{tabular}

\caption{Reflective-marker frequency and response length when the weight of the shared region is increased while the norm of the $\lambda=0$ adapter is kept fixed. $\gamma$ is the coefficient on the added shared-region component, and All in region moves the entire update into the shared region. The table tests whether thinking increases as the region is scaled up. Values after $\pm$ are standard deviations across three seeds.}
\label{tab:10}
\end{table*}

\subsection{Student Features}\label{sec:sae}

We therefore check whether the results above arise only from the frequency of reflective words, by asking whether the student's internal thinking features lead to the same conclusions. In the three Llama student cells, thinking features turn off at the $\lambda$ where each cell's teacher stops thinking ($0.2 \to 0.4$ in pair A, $0.4 \to 0.6$ in pair B, and in two steps, $0.4 \to 0.6 \to 0.8$, in pair C), and they move with reflective-marker frequency at $r = 0.88$ in pair A and $r = 0.92$ in pair C. Pair B gives $r = -0.08$: the features turn off, but a few reflective markers remain on the short-answer side (Table~\ref{tab:11}). Even in $\lambda=0$ answers with all reflective words deleted, 89--94\% of the original feature activation remains. Removing the shared region turns the features down to 13--24\% of the original in all three cells, whereas removing the same energy from other directions leaves 92--106\%, and raising the shared region's weight under a fixed norm keeps them at 78--115\%, around the $\lambda=0$ level. Grafting at $\beta=2$ turns them on only in pairs A and C (0.46 and 0.64, against 0.10--0.13 when grafting from other directions). B/Llama stays at the $\lambda=1$ level after grafting, and grafting other directions raises reflective markers to as many as 58 per 1,000 words through repetition while the features stay unchanged.

Features that do not turn on when reflective words are inserted and persist when they are deleted move at the same $\lambda$ as reflective-marker frequency, so the changes measured above are not changes in reflective words alone. Scaling up does not push the features above the $\lambda=0$ level (the 1.15$\times$ in pair A is comparable to the seed-to-seed spread of $\pm 0.09$), so the conclusion of \S\ref{sec:amplify} that the region does not control the amount of thinking also holds for thinking features. In pair A, removing other directions or raising the region's weight lowered reflective markers to 45--68\% of the original while the features stayed unchanged; this looks like a change in wording rather than in thinking, and the word frequencies in Tables~\ref{tab:6} and~\ref{tab:10} may underestimate the thinking of edited adapters. That grafting did not turn on the features in pair B suggests that the added reflective markers in pair B were repeated phrases rather than thinking. The selected features capture not only reflection but more broadly the way of thinking in which a solution is worked out in words, and because a public SAE exists only for Llama, the other two students are not examined. In short, read through thinking features, the switch, removal and amplification give the same conclusions in all three pairs, and grafting does so in pairs A and C, so the results above do not arise only from the frequency of reflective words.

\begin{table*}[t]
\centering
\small
% Table 11 (wide: use table*), generated from overleaf/05_what_the_shared_region_carries.md
\resizebox{\linewidth}{!}{%
\begin{tabular}{lrrrrrrrr}
\toprule
Pair & $\lambda$ = 0 & $\lambda$ = 1 & Shared removed & Other removed & Graft shared ($\beta$ = 2) & Graft other ($\beta$ = 2) & All in region & $\lambda$ = 0, markers deleted \\
\midrule
A & 0.89 $\pm$ 0.04 & 0.09 $\pm$ 0.002 & 0.21 $\pm$ 0.03 & 0.94 $\pm$ 0.01 & 0.46 $\pm$ 0.06 & 0.13 $\pm$ 0.01 & 1.02 $\pm$ 0.09 & 0.79 $\pm$ 0.04 \\
B & 0.68 $\pm$ 0.02 & 0.20 $\pm$ 0.01 & 0.14 $\pm$ 0.03 & 0.63 $\pm$ 0.03 & 0.19 $\pm$ 0.01 & 0.19 $\pm$ 0.02 & 0.53 $\pm$ 0.03 & 0.64 $\pm$ 0.01 \\
C & 0.93 $\pm$ 0.05 & 0.09 $\pm$ 0.003 & 0.13 $\pm$ 0.01 & 0.92 $\pm$ 0.01 & 0.64 $\pm$ 0.07 & 0.10 $\pm$ 0.01 & 0.92 $\pm$ 0.07 & 0.86 $\pm$ 0.05 \\
\bottomrule
\end{tabular}}

\caption{Mean activation of thinking features when the Llama-3.2-1B student reads its own answers, where 1 is the level on the $\lambda=0$ teacher's solutions. Edit conditions are as in Tables~\ref{tab:6} and~\ref{tab:8}--\ref{tab:10}, and the last column is the value after reflective words are deleted from the $\lambda=0$ answers. The features turn off when the shared region is removed and stay unchanged when other directions are removed or the shared region is scaled up. Values after $\pm$ are standard deviations across three seeds.}
\label{tab:11}
\end{table*}
